\section{Deneysel Kurulum}
\subsection{Çalışma Ortamı}
Deneyler Google Colab ortamında GPU hızlandırma kullanılarak yürütülmüştür. Uygulama PyTorch tabanlı olarak gerçekleştirilmiş; veri artırma için Albumentations kullanılmıştır \cite{buslaev2020albumentations}.

\subsection{Eğitim Ayarları}
Model eğitiminde aşağıdaki hiperparametreler kullanılmıştır:
\begin{itemize}
    \item Görüntü boyutu: $512\times512$,
    \item Batch boyutu: 4,
    \item Epoch sayısı: 15,
    \item Optimizasyon: Adam,
    \item Öğrenme oranı: $1\times10^{-3}$,
    \item Rastgelelik tohumu (seed): 42.
\end{itemize}

\begin{table}[H]
\centering
\caption{Eğitim hiperparametreleri.}
\label{tab:hyperparams}
\begin{tabular}{ll}
\toprule
\textbf{Parametre} & \textbf{Değer} \\
\midrule
Girdi çözünürlüğü & $512\times512$ \\
Batch boyutu & 4 \\
Epoch sayısı & 15 \\
Optimizasyon & Adam \\
Öğrenme oranı & $1\times10^{-3}$ \\
Seed & 42 \\
\bottomrule
\end{tabular}
\end{table}

\subsection{Model Seçimi ve Checkpoint}
Eğitim sürecinde doğrulama kümesi metrikleri izlenmiş ve en iyi doğrulama performansını veren model checkpoint olarak kaydedilmiştir. Test sonuçları, bu en iyi model üzerinden raporlanmıştır.
